\section*{Overview}

Hopcroft-Karp is the maximum matching algorithm I use when a bipartite graph is large enough that the simple Kuhn DFS
starts to feel slow. The main upgrade is phase-based augmentation: instead of finding one augmenting path at a time, it
finds many vertex-disjoint shortest augmenting paths in the same BFS/DFS round.

\section*{When to Use It}

Use it when:

\begin{itemize}
  \item the graph is bipartite,
  \item the goal is a maximum cardinality matching,
  \item \(O(VE)\) augmenting-path code is too slow.
\end{itemize}

\section*{Core Idea}

Alternate between two steps:

\begin{itemize}
  \item BFS from all unmatched left-side vertices to compute distances in the alternating graph,
  \item DFS that only follows edges consistent with those shortest distances to find a maximal set of shortest augmenting
  paths.
\end{itemize}

\section*{Key Insight}

After one phase, all shortest augmenting paths disappear. Therefore the distance from free left vertices to free right
vertices strictly increases from phase to phase, which limits how many phases can occur.

\section*{Worked Problem}

\paragraph{Problem.}
There are workers on the left, jobs on the right, and an edge means the worker can do that job. Find the largest number
of worker-job assignments with no repeated worker or job.

\paragraph{Why Hopcroft-Karp fits.}
Bipartite matching is exactly the right abstraction. The phase idea matters when the graph is dense or when both sides
have around \(10^5\) vertices total.

\section*{Correctness Intuition}

BFS layers the alternating graph by shortest augmenting-path length. DFS then uses only those shortest layers, so every
augmentation in the phase has minimum length. Because the chosen paths are vertex-disjoint, they can all be applied
simultaneously without conflict.

\section*{Complexity Analysis}

The complexity is \(O(E\sqrt{V})\).

\section*{Implementation}

The code stores the left partition explicitly and returns both the matching size and the matched partner arrays.

\section*{Common Pitfalls}

\begin{itemize}
  \item Forgetting which side the BFS starts from.
  \item Mixing 0-indexed and 1-indexed NIL conventions.
  \item Using Hopcroft-Karp for weighted matching; that is a different problem.
\end{itemize}

\section*{Variants / Extensions}

\begin{itemize}
  \item Minimum vertex cover from the final alternating-reachability state.
  \item Dinic on the flow network when you want one shared framework.
  \item Kuhn algorithm for smaller graphs or easier implementation.
\end{itemize}

\section*{Practice Problems}

\begin{itemize}
  \item Maximum bipartite matching.
  \item Assignment feasibility on an unweighted bipartite graph.
  \item Derive minimum vertex cover after computing the matching.
\end{itemize}

\section*{References}

\begin{itemize}
  \item \href{https://cp-algorithms.com/graph/kuhn_maximum_bipartite_matching.html}{cp-algorithms: Bipartite Matching}.
  \item \href{https://usaco.guide/plat/flow?lang=cpp}{USACO Guide: Flow and Matching}.
  \item \href{https://codeforces.com/catalog}{Codeforces Catalog}.
\end{itemize}
